\documentclass[11pt]{article}

\usepackage[margin=0.9in]{geometry}
\usepackage{amsmath, amssymb, mathtools}
\usepackage{hyperref}
\usepackage{enumitem}
\usepackage{parskip}
\usepackage{array}

\hypersetup{
    colorlinks=true,
    linkcolor=blue,
    urlcolor=blue
}

\title{CEC 315 --- Exam 3 Quick-Reference Notes\\\large Lectures 16--23}
\author{CEC 315 --- Signals and Systems (Spring 2026)}
\date{}

\newcommand{\Lap}{\overset{\mathcal{L}}{\longleftrightarrow}}
\newcommand{\Zed}{\overset{\mathcal{Z}}{\longleftrightarrow}}
\renewcommand{\Re}{\operatorname{Re}}

\begin{document}
\maketitle
\tableofcontents
\newpage

\section{Formula Tables}

\subsection{Laplace Transform Pairs (Causal)}
\begin{center}
\begin{tabular}{|l|l|l|}
\hline
$x(t)$ & $X(s)$ & ROC\\
\hline
$\delta(t)$ & $1$ & all $s$\\
$\delta(t-t_0)$ & $e^{-st_0}$ & all $s$\\
$u(t)$ & $1/s$ & $\Re\{s\}>0$\\
$tu(t)$ & $1/s^2$ & $\Re\{s\}>0$\\
$\dfrac{t^{n-1}}{(n-1)!}u(t)$ & $1/s^n$ & $\Re\{s\}>0$\\
$e^{-at}u(t)$ & $1/(s+a)$ & $\Re\{s\}>-\Re\{a\}$\\
$te^{-at}u(t)$ & $1/(s+a)^2$ & $\Re\{s\}>-\Re\{a\}$\\
$\cos(\omega_0 t)u(t)$ & $s/(s^2+\omega_0^2)$ & $\Re\{s\}>0$\\
$\sin(\omega_0 t)u(t)$ & $\omega_0/(s^2+\omega_0^2)$ & $\Re\{s\}>0$\\
$e^{-at}\cos(\omega_d t)u(t)$ & $(s+a)/[(s+a)^2+\omega_d^2]$ & $\Re\{s\}>-a$\\
$e^{-at}\sin(\omega_d t)u(t)$ & $\omega_d/[(s+a)^2+\omega_d^2]$ & $\Re\{s\}>-a$\\
$-e^{-at}u(-t)$ & $1/(s+a)$ & $\Re\{s\}<-\Re\{a\}$\\
\hline
\end{tabular}
\end{center}

\subsection{Laplace Properties}
\begin{center}
\begin{tabular}{|l|l|l|}
\hline
Property & Time & $s$-Domain\\
\hline
Linearity & $ax+by$ & $aX+bY$\\
Time shift & $x(t-t_0)$ & $e^{-st_0}X(s)$\\
$s$-shift & $e^{s_0 t}x(t)$ & $X(s-s_0)$\\
Scaling ($a>0$) & $x(at)$ & $X(s/a)/a$\\
Convolution & $x*y$ & $XY$\\
Diff. (bil.) & $dx/dt$ & $sX$\\
Diff. (unil.) & $dx/dt$ & $sX - x(0^-)$\\
2nd diff.\ (unil.) & $d^2x/dt^2$ & $s^2 X - sx(0^-)-x'(0^-)$\\
Integration & $\int_{-\infty}^t x\,d\tau$ & $X/s$\\
$s$-diff. & $-tx(t)$ & $dX/ds$\\
IVT & $x(0^+)$ & $\lim_{s\to\infty}sX(s)$\\
FVT & $x(\infty)$ & $\lim_{s\to 0}sX(s)$\\
\hline
\end{tabular}
\end{center}

\subsection{Z-Transform Pairs}
\begin{center}
\begin{tabular}{|l|l|l|}
\hline
$x[n]$ & $X(z)$ & ROC\\
\hline
$\delta[n]$ & $1$ & all $z$\\
$\delta[n-k]$ & $z^{-k}$ & $|z|>0$\\
$u[n]$ & $1/(1-z^{-1})$ & $|z|>1$\\
$a^n u[n]$ & $1/(1-az^{-1})$ & $|z|>|a|$\\
$-a^n u[-n-1]$ & $1/(1-az^{-1})$ & $|z|<|a|$\\
$na^n u[n]$ & $az^{-1}/(1-az^{-1})^2$ & $|z|>|a|$\\
$(n+1)a^n u[n]$ & $1/(1-az^{-1})^2$ & $|z|>|a|$\\
$r^n\cos(\omega_0 n)u[n]$ & $\dfrac{1-r\cos\omega_0\, z^{-1}}{1-2r\cos\omega_0\, z^{-1}+r^2 z^{-2}}$ & $|z|>r$\\
$r^n\sin(\omega_0 n)u[n]$ & $\dfrac{r\sin\omega_0\, z^{-1}}{1-2r\cos\omega_0\, z^{-1}+r^2 z^{-2}}$ & $|z|>r$\\
\hline
\end{tabular}
\end{center}

\subsection{Z-Transform Properties}
\begin{center}
\begin{tabular}{|l|l|l|}
\hline
Property & Sequence & z-Domain\\
\hline
Linearity & $ax_1+bx_2$ & $aX_1+bX_2$\\
Time shift & $x[n-n_0]$ & $z^{-n_0}X$\\
z-scaling & $z_0^n x[n]$ & $X(z/z_0)$\\
Time reversal & $x[-n]$ & $X(z^{-1})$\\
Convolution & $x_1*x_2$ & $X_1 X_2$\\
$z$-diff. & $nx[n]$ & $-z\,dX/dz$\\
First diff. & $x[n]-x[n-1]$ & $(1-z^{-1})X$\\
Accumulation & $\sum_{k\le n}x[k]$ & $X/(1-z^{-1})$\\
Unilateral shift & $x[n-1]$ & $z^{-1}X + x[-1]$\\
Unil. shift-2 & $x[n-2]$ & $z^{-2}X + x[-1]z^{-1}+x[-2]$\\
IVT & $x[0]$ & $\lim_{z\to\infty}X(z)$\\
FVT & $x[\infty]$ & $\lim_{z\to 1}(1-z^{-1})X(z)$\\
\hline
\end{tabular}
\end{center}

\section{Per-Lecture Reference}

\subsection{Lecture 16 --- Laplace Transform and ROC}
\begin{itemize}[leftmargin=*]
  \item PDF: \texttt{all\_lectures/cec315-lctr16-laplace-transform-roc.pdf}
  \item Summary: \texttt{lecture\_summaries/lctr16\_laplace\_transform\_roc.md}
  \item Textbook: Oppenheim \S 9.0--9.2 (pp.~654--670).
\end{itemize}
\textbf{Key ideas:}
\begin{itemize}[leftmargin=*]
  \item $X(s) = \int_{-\infty}^\infty x(t)e^{-st}\,dt$, with $s = \sigma + j\omega$.
  \item ROC is a vertical strip; cannot contain poles.
  \item Right-sided $\Rightarrow$ ROC right of rightmost pole; left-sided $\Rightarrow$ left of leftmost pole.
  \item Fourier transform exists iff $j\omega$-axis $\subset$ ROC.
\end{itemize}

\subsection{Lecture 17 --- Inverse Laplace and Properties}
\begin{itemize}[leftmargin=*]
  \item PDF: \texttt{all\_lectures/cec315-lctr17-inverse-laplace-properties.pdf}
  \item Summary: \texttt{lecture\_summaries/lctr17\_inverse\_laplace\_properties.md}
  \item Textbook: \S 9.3--9.6 (pp.~670--692).
\end{itemize}
\textbf{Key ideas:}
\begin{itemize}[leftmargin=*]
  \item Practical inversion = partial fractions + table lookup.
  \item ROC determines direction of each term.
  \item Complex conjugate poles $\Rightarrow$ complete the square, use $\cos/\sin$ pairs.
  \item IVT/FVT give $x(0^+)$ and $x(\infty)$ directly.
  \item Convolution $\leftrightarrow$ multiplication; differentiation $\leftrightarrow$ multiply by $s$.
\end{itemize}

\subsection{Lecture 18 --- System Analysis via Unilateral Laplace}
\begin{itemize}[leftmargin=*]
  \item PDF: \texttt{all\_lectures/cec315-lctr18-system-analysis-unilateral-laplace.pdf}
  \item Summary: \texttt{lecture\_summaries/lctr18\_system\_analysis\_unilateral\_laplace.md}
  \item Textbook: \S 9.7--9.9 (pp.~693--720).
\end{itemize}
\textbf{Key ideas:}
\begin{itemize}[leftmargin=*]
  \item $H(s) = Y/X$ by $d^k/dt^k \to s^k$ (zero ICs).
  \item Causal + stable $\Leftrightarrow$ all poles in open LHP.
  \item Block diagrams: series $\times$, parallel $+$, feedback $H/(1+GH)$.
  \item Unilateral: $\mathcal{L}_u\{y'\} = sY - y(0^-)$; $\mathcal{L}_u\{y''\} = s^2Y - sy(0^-) - y'(0^-)$.
  \item Total response = ZSR + ZIR.
\end{itemize}

\subsection{Lecture 19 --- Z-Transform and ROC}
\begin{itemize}[leftmargin=*]
  \item PDF: \texttt{all\_lectures/cec315-lctr19-z-transform-roc.pdf}
  \item Summary: \texttt{lecture\_summaries/lctr19\_z\_transform\_roc.md}
  \item Textbook: \S 10.0--10.2 (pp.~741--757).
\end{itemize}
\textbf{Key ideas:}
\begin{itemize}[leftmargin=*]
  \item $X(z) = \sum_n x[n] z^{-n}$, $z = re^{j\omega}$.
  \item Unit circle $|z|=1$ plays the role of the $j\omega$-axis.
  \item ROC is annular ring; right-sided exterior, left-sided interior.
  \item DTFT exists iff unit circle $\subset$ ROC.
\end{itemize}

\subsection{Lecture 20 --- Inverse Z-Transform and Properties}
\begin{itemize}[leftmargin=*]
  \item PDF: \texttt{all\_lectures/cec315-lctr20-inverse-z-transform-properties.pdf}
  \item Summary: \texttt{lecture\_summaries/lctr20\_inverse\_z\_transform\_properties.md}
  \item Textbook: \S 10.3--10.6 (pp.~757--774).
\end{itemize}
\textbf{Key ideas:}
\begin{itemize}[leftmargin=*]
  \item Partial fractions in $z^{-1}$, not $z$.
  \item ROC outside pole $\Rightarrow$ right-sided; inside $\Rightarrow$ left-sided.
  \item Complex poles at $re^{\pm j\omega_0}$ $\Rightarrow$ damped sinusoid.
  \item $z^{-1}$ is the unit-delay operator (DSP).
\end{itemize}

\subsection{Lecture 21 --- System Analysis via Unilateral Z}
\begin{itemize}[leftmargin=*]
  \item PDF: \texttt{all\_lectures/cec315-lctr21-system-analysis-unilateral-z.pdf}
  \item Summary: \texttt{lecture\_summaries/lctr21\_system\_analysis\_unilateral\_z.md}
  \item Textbook: \S 10.7--10.9 (pp.~774--796).
\end{itemize}
\textbf{Key ideas:}
\begin{itemize}[leftmargin=*]
  \item $H(z) = Y/X$ from difference equation by $y[n-k]\to z^{-k}Y$.
  \item Causal + stable $\Leftrightarrow$ all poles inside unit circle.
  \item Unilateral shift: $x[n-1] \to z^{-1}X + x[-1]$ (plus, not minus!).
  \item Total = ZSR + ZIR.
\end{itemize}

\subsection{Lecture 22 --- Sampling}
\begin{itemize}[leftmargin=*]
  \item PDF: \texttt{all\_lectures/cec315-lctr22-sampling.pdf}
  \item Summary: \texttt{lecture\_summaries/lctr22\_sampling.md}
  \item Textbook: Ch.~7 (pp.~514--555).
\end{itemize}
\textbf{Key ideas:}
\begin{itemize}[leftmargin=*]
  \item $X_p(j\omega) = \frac{1}{T}\sum_k X(j(\omega - k\omega_s))$ (replicated spectrum).
  \item Sampling theorem: $\omega_s > 2\omega_M$ (strict!).
  \item Ideal reconstruction: LPF (gain $T$) + sinc interpolation.
  \item Aliasing irreversible; use anti-aliasing filter before sampling.
  \item $\Omega = \omega T$ (DT/CT frequency relation).
\end{itemize}

\subsection{Lecture 23 --- Linear Feedback Systems}
\begin{itemize}[leftmargin=*]
  \item PDF: \texttt{all\_lectures/cec315-lctr23-linear-feedback-systems.pdf}
  \item Summary: \texttt{lecture\_summaries/lctr23\_linear\_feedback\_systems.md}
  \item Textbook: Ch.~11 \S 11.0--11.5 (pp.~816--866).
\end{itemize}
\textbf{Key ideas:}
\begin{itemize}[leftmargin=*]
  \item $Q(s) = H/(1+GH)$; characteristic equation $1+GH = 0$.
  \item Block diagram: series $\times$, parallel $+$, feedback $H/(1+GH)$.
  \item Nyquist: no encirclement of $-1/K$ $\Rightarrow$ stable (assuming OL stable).
  \item GM: dB below 0 at $\omega_1$ ($\angle GH = -180^\circ$); PM: degrees above $-180^\circ$ at $\omega_2$ ($|GH|=0$ dB).
  \item Both positive $\Rightarrow$ stable. $\tau_{\max} = \text{PM (rad)}/\omega_2$.
\end{itemize}

\section{Pitfalls and Tricky Edge Cases}

\subsection{Transform Pitfalls}
\begin{enumerate}[leftmargin=*]
  \item \textbf{Always state the ROC.} $1/(s+3)$ is ambiguous without it.
  \item ROC (not pole sign) determines direction in partial-fraction inversion.
  \item Repeated poles of order $n$ require $n$ terms $B_k/(s-p)^k$.
  \item Complex poles: complete the square; don't split into complex partial fractions.
  \item Poles of $1/(1-az^{-1})$ are at $z=a$, not $1/a$.
  \item Z partial fractions: work in $z^{-1}$, not $z$.
\end{enumerate}

\subsection{Unilateral Sign Conventions}
\begin{itemize}[leftmargin=*]
  \item Laplace: $\mathcal{L}_u\{y'\} = sY - y(0^-)$ (\textbf{minus}).
  \item Z: $\mathcal{Z}_u\{y[n-1]\} = z^{-1}Y + y[-1]$ (\textbf{plus}).
\end{itemize}

\subsection{Stability Confusions}
\begin{itemize}[leftmargin=*]
  \item CT: open LHP. DT: interior of unit circle.
  \item $z=-0.9$ is stable (magnitude $0.9$), though on negative real axis.
  \item Marginal stability (poles on $j\omega$-axis / unit circle) is \emph{not} BIBO stable.
  \item RHP zeros do not cause instability; only poles matter.
\end{itemize}

\subsection{Final-Value Theorem Traps}
Only valid if $sX(s)$ (or $(1-z^{-1})X(z)$) has all poles strictly in open LHP (inside unit circle). If the signal oscillates or grows, FVT gives meaningless result.

\subsection{Sampling Traps}
\begin{itemize}[leftmargin=*]
  \item $\omega_s = 2\omega_M$ is \emph{not} sufficient.
  \item Mixing Hz and rad/s causes $2\pi$ errors.
  \item Aliasing is irreversible; anti-aliasing must happen before sampling.
  \item Squaring a signal doubles the bandwidth.
  \item Reconstruction LPF has gain $T$ (cancels $1/T$ scaling).
  \item Nyquist rate of $x(t-a)$ same as $x$; Nyquist rate of $x(2t)$ is doubled.
\end{itemize}

\subsection{Feedback / Margin Traps}
\begin{itemize}[leftmargin=*]
  \item Negative feedback $\to 1+GH$; positive $\to 1-GH$.
  \item Open-loop poles (of $GH$) $\neq$ closed-loop poles (roots of $1+KGH=0$).
  \item GM on magnitude plot at $\omega_1$ where $\angle GH = -180^\circ$.
  \item PM on phase plot at $\omega_2$ where $|GH|=0$ dB.
  \item Feedback does not always stabilize; too much gain can destabilize.
  \item First-order plants $\Rightarrow$ infinite gain margin.
  \item $\tau_{\max} = \text{PM (rad)}/\omega_2$.
\end{itemize}

\subsection{ROC Reminders}
\begin{itemize}[leftmargin=*]
  \item ROC never contains poles.
  \item Right-sided: right of rightmost pole (CT) / outside outermost pole (DT).
  \item Left-sided: left of leftmost pole (CT) / inside innermost pole (DT).
  \item Two-sided: strip/ring between consecutive poles; may be empty.
  \item Finite-duration: entire plane (except possibly $z=0,\infty$).
\end{itemize}

\end{document}
